\honest{Hug the Query}{Eliezer Yudkowsky}{2007}

Evidence close to a question, I say, beats evidence further away. The Wright brothers say their plane will fly. Compare their credentials with Lord Kelvin's, and Kelvin wins. Check both sides' calculations, and credentials matter less; Kelvin ``probably doesn't have any apart from his own incredulity.'' Watch the plane fly, and the calculations become moot ``for many purposes.''

I cite no statement by Kelvin about the Wrights. And it was the Wrights' calculations that failed. Their 1901 glider gave about a third of the lift they had calculated, because they used a coefficient that had been accepted for more than a century and was wrong. They found the error by building a wind tunnel, not by checking their arithmetic.

``It's a theorem of these causal graphs that you can never get more information from distant nodes, than from strictly closer nodes that screen off the distant ones.'' That is true when the closer evidence really does screen off the distant. It says nothing about reading the closer evidence wrongly.

A writing teacher, Jerry Cleaver, said that what does you in is ``overlooking the basics.''

It is better to argue physics than credentials, I say, and better to argue physics than rationality. Accusing someone of ``Bias \#182'' cannot settle a question of fact. This is good advice.

Then the concession. You cannot always check the calculations: you may lack the background, the information may be private, or there may be no time. ``There's a sadly large number of times when it's worthwhile to judge the speaker's rationality.'' For those times I offer ``a hollow feeling in your heart.'' That is the ordinary situation of anyone judging a question in medicine or climate science, and a hollow feeling is not a method.

In a footnote to that concession, I point to my own post defending molecular nanotechnology, a claim no experiment had yet tested.

``Whenever you can, dance as near to the original question as possible---press yourself up against it---get close enough to hug the query!''
